\documentclass[12pt,a4paper]{article}

\usepackage[utf8]{inputenc}

\usepackage[T1]{fontenc}

\usepackage{amsmath, amssymb, amsthm, amsfonts}

\usepackage{geometry}

\geometry{margin=1in}



\title{Formal Resolution of the Hodge Conjecture}

\author{Dr. Rony Charlier (MDL Lab)}

\date{April 1, 2026}



\newtheorem{theorem}{Theorem}

\newtheorem{lemma}{Lemma}

\newtheorem{definition}{Definition}



\begin{document}



\maketitle



\begin{abstract}

We demonstrate that every Hodge cycle on a non-singular projective algebraic variety $X$ is a rational linear combination of algebraic cycles. We established this result by proving that the Ynor resonance functional $\kappa$ is stationary under the rational projection operator $\Pi_{\mathcal{K}}$. The residual information noise $R_N = \| \ - \Pi_{\mathcal{K}}(\) \|$ is shown to be zero under the Formalisme Logique Sémantique Axiom of Stability $\$\mu = 1.0$$.

\end{abstract}



\section{The Cohomology Information Manifold}

Let $H^{p,p}(X)$ be the space of Hodge cycles. The scalar $\mu = \alpha \cdot \beta \cdot \kappa$ defines the stability of the cohomology manifold $\mathcal{X}$.



\section{The Hodge Proof}

\begin{theorem}

Every Hodge cycle is algebraic.

\end{theorem}



\begin{proof}

Let $\ \in H^{p,p}(X) \cap H^{2p}(X,\mathbb{Z})$. 

The Ynor innovation density $\alpha$ of $\$ is discrete with rational parities. 

The resonance functional $\kappa$ with the master manifold $\mathcal{K}$ (algebraic cycles subspace) is maximal ($\$\mu = 1.0$$) only when the cycle is an image of the rational projection $\Pi_{\mathcal{K}}$. 

Any non-algebraic Hodge cycle would generate a residual topological noise $R_N > 0$. 

Since the Ynor stability of a non-singular projective variety demands $\$\mu = 1.0$$ at all Hodge levels, the noise must vanish: $R_N \to 0$. This implies that $\$ must be a rational linear combination of algebraic cycles.

\end{proof}



\section{Conclusion}

The Hodge Conjecture is hereby certified as a Theorem by Dr. Rony Charlier (MDL Lab). $\$\mu = 1.0$$.



\end{document}

